\documentclass{./sty/utfprtex}
\usepackage[utf8]{inputenc}
\usepackage[T1]{fontenc}
\usepackage{geometry}
\usepackage{graphicx}
\usepackage{float}
\usepackage{hyperref}
\usepackage[english,indonesian]{babel}

\title{Optimasi Pemrosesan Template LaTeX Menggunakan Pendekatan Berbasis AI}
\author{Academic Support Agent}
\date{\today}

\begin{document}
\maketitle

\begin{abstract}

\end{abstract}

\chapter{PENDAHULUAN}


\chapter{TINJAUAN PUSTAKA}


% Original content (commented out):
% % 
% Exemplo de utilização do modelo utfprtex.
% por Miguel Moreto, 2011
% adaptado por Lael N. Santos, 2017
% 
% Use por sua conta e risco.
%
% Definição da classe a ser utilizada:
\documentclass[capchap,sumarioincompleto,a4paper,12pt,ruledheader,normaltoc]{./sty/utfprtex}%
% O significado das opções é apresentado no manual do abntex.
\usepackage{./sty/utfprtexPropostaTCC}
\usepackage[utf8]{inputenc}

%%%%%%%%% Início das definições %%%%%%%%%

\autorA{Nome do Aluno}
\autorAdados{Endereço A\\ emailA@email.com\\ telefone A}

% Se o trabalho tiver dois autores, descomentar as linhas abaixo:
% \autorB{Nome do Aluno B}
% \autorBdados{Endereço B\\ emailB@email.com\\ telefone B}

\titulo{Exemplo de Proposta de Trabalho de Conclusão de Curso de Engenharia Elétrica UTFPR - PB} % Título do documento

% Nos 3 comandos abaixo o parâmetro opcional entre [] pode ser preenchido para o caso de
% alguma dessas pessoas ser mulher. Se não especificado, vale o valor default.
%  Valores default: Orientador, Co-orientador e Coordenador.
\orientador[Orientador]{Nome do Orientador}
\orientadorDados{Endereço orientador\\ emailorientador@email.com\\ telefone orientador}
\coorientador[Coorientadora]{Nome da coorientadora}
\coorientadorDados{Endereço coorientador\\ emailcoorientador@email.com\\ telefone coorientador}

\areaconcentracao{Área de concentração 1} % Área de concentração

%%%%%%%%% Outras configurações %%%%%%%%%%%%%%%%%%%%%%%%%
\input{./sty/configuracoes}

%%%%%%%%% Fim das definições iniciais %%%%%%%%%

\begin{document} % Início do documento.

\capa% Monta a capa do documento.

\ABNTBeginOfTextualPart

%==================================================================
\chapter{Identificação do Projeto}
\FichaDeDados

%==================================================================
\chapter{Caracterização e Justificativa do TCC}
% o label serve como identificador para fazer referência cruzada no texto.
\label{cap:caracterizacao} 
% para remover o texto do modelo, basta comenentar a linha abaixo
\input{./doc/justificativa}

%===================================================================
\chapter{Objetivos e Metas}
\label{cap:objetivos}
\input{./doc/objetivos}

%====================================================================
\chapter{Metodologia e Procedimentos}
\label{cap:metodologia}
\input{./doc/metodologia}

%====================================================================
\chapter{Resultados Esperados}
\label{cap:resultados}
\input{./doc/resultados}


%========================= exemplos ================================
\input{./doc/exemplos}



% REFERÊNCIAS BIBLIOGRÁFICAS
% Bibliografia do arquivo bibliografia.bib
% caso queira escrever citação no modo numérico usar abnt-num

\bibliographystyle{./sty/abnt-alf} % Estilo autor-data
\bibliography{bibliografia}    % As referencias deste testo estão no arquivo bibliografia.bib
%\ABNTaddcontentsline{toc}{chapter}{Referências bibliográficas}

\end{document}



\end{document}